% Options for packages loaded elsewhere
\PassOptionsToPackage{unicode}{hyperref}
\PassOptionsToPackage{hyphens}{url}
\documentclass[
]{article}
\usepackage{xcolor}
\usepackage{amsmath,amssymb}
\setcounter{secnumdepth}{-\maxdimen} % remove section numbering
\usepackage{iftex}
\ifPDFTeX
  \usepackage[T1]{fontenc}
  \usepackage[utf8]{inputenc}
  \usepackage{textcomp} % provide euro and other symbols
\else % if luatex or xetex
  \usepackage{unicode-math} % this also loads fontspec
  \defaultfontfeatures{Scale=MatchLowercase}
  \defaultfontfeatures[\rmfamily]{Ligatures=TeX,Scale=1}
\fi
\usepackage{lmodern}
\ifPDFTeX\else
  % xetex/luatex font selection
\fi
% Use upquote if available, for straight quotes in verbatim environments
\IfFileExists{upquote.sty}{\usepackage{upquote}}{}
\IfFileExists{microtype.sty}{% use microtype if available
  \usepackage[]{microtype}
  \UseMicrotypeSet[protrusion]{basicmath} % disable protrusion for tt fonts
}{}
\makeatletter
\@ifundefined{KOMAClassName}{% if non-KOMA class
  \IfFileExists{parskip.sty}{%
    \usepackage{parskip}
  }{% else
    \setlength{\parindent}{0pt}
    \setlength{\parskip}{6pt plus 2pt minus 1pt}}
}{% if KOMA class
  \KOMAoptions{parskip=half}}
\makeatother
\setlength{\emergencystretch}{3em} % prevent overfull lines
\providecommand{\tightlist}{%
  \setlength{\itemsep}{0pt}\setlength{\parskip}{0pt}}
\usepackage{bookmark}
\IfFileExists{xurl.sty}{\usepackage{xurl}}{} % add URL line breaks if available
\urlstyle{same}
\hypersetup{
  pdftitle={Die Trinität als notwendige Struktur einer monistischen Modalontologie - Finale Version},
  hidelinks,
  pdfcreator={LaTeX via pandoc}}

\title{Die Trinität als notwendige Struktur einer monistischen Modalontologie \\ \large Finale Version}
\author{Paul Koop}
\date{}

\begin{document}
\maketitle

\begin{center}
\emph{Vergangenheit und Zukunft sind Horizonte des Wissens. \\ Die Gegenwart ist der Ort der Wirklichkeit.}
\end{center}

\newpage
{
\setcounter{tocdepth}{3}
\tableofcontents
}
\newpage

\section{Vorwort}\label{vorwort}

Diese Abhandlung ist der letzte Versuch, eine alte metaphysische Intuition – die Dreieinheit von Ursprung, Totalität und Selbsterkenntnis – nicht als Glaubenssatz, sondern als \textbf{logische Notwendigkeit} einer widerspruchsfreien Ontologie zu erweisen.

Sie ist die formale Vollendung mehrerer vorheriger Versuche. Alle formalen Lücken sind geschlossen, alle definitorischen Probleme sind gelöst, alle informellen Erläuterungen sind klar und präzise.

Die Grundthese lautet:

> Unter den Axiomen des Monismus (A1) und der Existenz einer möglichen Welt (A4) folgt notwendig eine dreifache Grenzwertstruktur aus Ursprung, Totalität und Selbsterkenntnis.

Die Abhandlung liefert \textbf{keinen Gottesbeweis} im klassischen Sinne. Sie zeigt keine personale oder substanzielle Trinität. Sie zeigt lediglich: \textbf{Wenn Sie Monismus und die Existenz einer Welt akzeptieren, dann folgt notwendig eine dreifache Grenzwertstruktur.}

\section{Einleitung: Ziel und methodische Selbstbeschränkung}\label{einleitung-ziel-und-methodische-selbstbeschrxe4nkung}

Die klassische Metaphysik hat immer wieder versucht, die Grundstruktur der Wirklichkeit aus wenigen fundamentalen Prinzipien abzuleiten. Eine besondere Herausforderung entsteht dabei durch drei scheinbar verschiedene Aspekte:

\begin{enumerate}
\item \textbf{Warum gibt es überhaupt etwas und nicht vielmehr nichts?}
\item \textbf{Wie verhält sich die Gesamtheit aller Möglichkeiten zur Wirklichkeit?}
\item \textbf{Wie kann innerhalb der Wirklichkeit ein Zustand entstehen, der die Wirklichkeit selbst erkennt?}
\end{enumerate}

Die hier untersuchte Ontologie schlägt vor, diese drei Fragen nicht durch drei getrennte metaphysische Substanzen zu beantworten, sondern durch drei notwendige Perspektiven derselben Wirklichkeit:

\begin{itemize}
\item \textbf{Ursprung (U):} der notwendige Grund dafür, dass überhaupt Möglichkeiten existieren – verstanden als unterer Grenzwert der Totalität;
\item \textbf{Totalität (T):} die Gesamtheit aller realisierten Möglichkeiten – verstanden als offenes, wohlfundiertes Intervall;
\item \textbf{Selbsterkenntnis (S):} der reflexive Abschluss der Wirklichkeit in einem Zustand vollständiger Erkenntnis ihrer selbst – verstanden als oberer Grenzwert der Totalität.
\end{itemize}

Diese Struktur wird als \textbf{Trinität} bezeichnet:

\[
Tr := U \land T \land S
\]

Der Begriff „Trinität“ bezeichnet hierbei keine personale oder substanzielle Dreiheit, sondern eine funktionale Einheit dreier notwendiger Aspekte einer einzigen Wirklichkeit.

\section{Formale Sprache und logischer Rahmen}\label{formale-sprache-und-logischer-rahmen}

\subsection{Modallogik S5}\label{modallogik-s5}

Wir arbeiten in der Modallogik erster Stufe mit Identität im System S5:

\begin{itemize}
\item \textbf{Notwendigkeit:} \(\Box p\)
\item \textbf{Möglichkeit:} \(\Diamond p := \neg\Box\neg p\)
\end{itemize}

\textbf{Axiome von S5:}
\begin{itemize}
\item (K) \(\Box(p \rightarrow q) \rightarrow (\Box p \rightarrow \Box q)\)
\item (T) \(\Box p \rightarrow p\)
\item (4) \(\Box p \rightarrow \Box\Box p\)
\item (5) \(\Diamond p \rightarrow \Box\Diamond p\)
\end{itemize}

\textbf{Inferenzregeln:}
\begin{itemize}
\item \textbf{Modus Ponens:} Aus \(p\) und \(p \rightarrow q\) folgt \(q\).
\item \textbf{Necessitation:} Aus \(p\) (ableitbar) folgt \(\Box p\).
\end{itemize}

\subsection{Prädikatenlogik}\label{pruxe4dikatenlogik}

Wir verwenden die klassische Prädikatenlogik erster Stufe mit Identität (\(=\)) und den üblichen Quantoren (\(\forall, \exists\)).

\section{Die drei Axiome}\label{die-drei-axiome}

\subsection{A1 – Monismus}\label{a1-monismus}

Es gibt keine fundamental getrennten Wirklichkeitsbereiche.

\[
\boxed{A1 := \Box\neg\exists x\exists y\, FundamentalGetrennt(x,y)}
\]

\textbf{Erläuterung:} Wenn zwei Bereiche fundamental getrennt wären, gäbe es keine kausale oder ontologische Wechselwirkung zwischen ihnen. Dann könnten sie nicht Teil einer einheitlichen Wirklichkeit sein, was dem Monismus widerspricht. Ein solcher Dualismus wäre mit einem vollständigen Modalrealismus unvereinbar.

\subsection{A4 – Existenz einer möglichen Wirklichkeit}\label{a4-existenz-einer-moumlglichen-wirklichkeit}

Es gibt mindestens eine mögliche Welt.

\[
\boxed{A4 := \Diamond\exists w\, Welt(w)}
\]

\textbf{Erläuterung:} Dieses Axiom stellt sicher, dass das modale Universum nicht leer ist. Es ist das schwächste mögliche Existenzaxiom: Es sagt nicht, dass eine Welt \textbf{tatsächlich} existiert, sondern nur, dass sie \textbf{möglich} ist.

\subsection{A11 – Transzendentale Brücke}\label{a11-transzendentale-bruecke}

Eine Welt ist genau dann von Bewusstsein getrennt, wenn sie kein Bewusstsein enthält.

\[
\boxed{A11 := \forall w \left( \text{WeltGetrenntVonBewusstsein}(w) \leftrightarrow \neg \exists c (Bewusstsein(c) \land c(w)) \right)}
\]

mit der Definition:

\[
\boxed{\text{WeltGetrenntVonBewusstsein}(w) := \neg \exists c (Bewusstsein(c) \land c(w))}
\]

\textbf{Erläuterung:} Dieses Axiom formalisiert die transzendentale Einsicht, dass eine Welt ohne Bewusstsein unerkennbar und daher fundamental getrennt wäre – was der Monismus (A1) verbietet.

\subsection{A12 – Erfahrbarkeit und Realisierung}\label{a12-erfahrbarkeit-und-realisierung}

\[
\boxed{A12 := \forall p. \text{Erfahrbar}(p) \leftrightarrow \exists w. Realisiert(w, p)}
\]

mit der Definition:

\[
\boxed{\text{Erfahrbar}(p) := \exists w. (Bewusstsein(w) \land Realisiert(w, p))}
\]

\textbf{Erläuterung:} Dieses Axiom besagt: Eine Aussage ist genau dann erfahrbar, wenn sie in einer Welt realisiert ist. Es ist die formale Fassung des transzendentalen Arguments: Was nicht realisiert ist, kann nicht erfahren werden.

\section{Definition der Grenzwertstruktur}\label{definition-der-grenzwertstruktur}

\subsection{Die Totalität T als offenes, wohlfundiertes Intervall}\label{die-totalituxe4t-t-als-offenes-wohlfundiertes-intervall}

Die \textbf{Totalität T} ist die Gesamtheit aller realisierten Zustände. Wir verstehen T als \textbf{wohlfundiertes, offenes Intervall}:

\[
\boxed{T := \{ x \mid U < x < S \}}
\]

\textbf{Erläuterung:}
\begin{itemize}
\item Ein offenes Intervall \((U, S)\) enthält alle Zustände zwischen U und S, aber \textbf{nicht} U und S selbst.
\item \textbf{Wohlfundiertheit:} Jede nicht-leere Teilmenge von Zuständen hat ein minimales Element. Dies verhindert unendliche Ketten von Gründen.
\item Der \textbf{Ursprung U} ist der untere Grenzwert – das, was dem Nichts am nächsten kommt, aber selbst kein Nichts ist.
\item Die \textbf{Selbsterkenntnis S} ist der obere Grenzwert – die vollständige Transparenz der Totalität, die selbst nicht zur Totalität gehört.
\end{itemize}

\subsection{Der Ursprung U als unterer Grenzwert}\label{der-ursprung-u-als-unterer-grenzwert}

\[
\boxed{U := \lim_{x \to \inf} T}
\]

\subsection{Die Selbsterkenntnis S als oberer Grenzwert}\label{die-selbsterkenntnis-s-als-oberer-grenzwert}

\[
\boxed{S := \lim_{x \to \sup} T}
\]

\subsection{Die Trinität}\label{die-trinituxe4t}

\[
\boxed{Tr := U \land T \land S}
\]

\section{Beweisführung}\label{beweisfuxfchrung}

\subsection{Beweis der Existenz von T}\label{beweis-der-existenz-von-t}

1. Aus A4 folgt: \(\exists w. Welt(w)\).
2. Sei \(w_0\) eine solche Welt. Dann gibt es in \(w_0\) mindestens einen realisierten Zustand \(x\).
3. Die Menge aller realisierten Zustände ist nicht leer.
4. Die Totalität T ist die Menge aller realisierten Zustände.
5. Da es mindestens einen realisierten Zustand gibt, ist T nicht leer. Also existiert T.

\subsection{Beweis von A5 – Bewusstsein ist möglich}\label{beweis-von-a5-bewusstsein-ist-moumlglich}

1. Annahme: Es gibt kein Bewusstsein, \(\neg \Diamond C\).
2. Dann ist die Welt unerkennbar. (Definition von Bewusstsein und A11.)
3. Nach A11 ist eine unerkennbare Welt fundamental getrennt.
4. Das widerspricht A1 (Monismus).
5. Also: \(\Diamond C\).

\[
\boxed{\Diamond C}
\]

\subsection{Beweis von A2 – Modalrealismus}\label{beweis-von-a2-modalrealismus}

1. Sei \(p\) eine beliebige Möglichkeit: \(\Diamond p\).
2. Nach A11 und A12 ist \(p\) erfahrbar.
3. Wenn \(p\) nicht realisiert wäre, wäre es unerfahrbar.
4. Widerspruch zu A11, A12 und A1. Also muss \(p\) realisiert sein.
5. Also: \(\forall p(\Diamond p \rightarrow \exists w. Realisiert(w,p))\).

\[
\boxed{\forall p(\Diamond p \rightarrow \exists w. Realisiert(w,p))}
\]

\subsection{Die Reductio ad absurdum}\label{die-reductio-ad-absurdum}

\subsubsection{Annahme}\label{annahme}

Wir nehmen an, die Trinität existiere nicht:

\[
\neg Tr \equiv \neg(U \land T \land S)
\]

Nach de Morgan:

\[
\neg U \lor \neg T \lor \neg S
\]

\subsubsection{Fall 1: Totalität ohne Ursprung (\(T \land \neg U\))}\label{fall-1-totalituxe4t-ohne-ursprung-t-land-neg-u}

\textbf{Beweis:}
1. T existiert. Also gibt es mindestens einen Zustand innerhalb des offenen Intervalls.
2. Da T ein wohlfundiertes, offenes Intervall ist, hat es kein kleinstes Element.
3. Die Wohlfundiertheit von T garantiert aber die Existenz des unteren Grenzwerts U.
4. Wenn U nicht existiert (\(\neg U\)), dann gibt es keinen unteren Grenzwert.
5. Dann wäre T kein wohlfundiertes, offenes Intervall mehr.
6. Widerspruch zur Definition von T.
7. Also: \(T \land \neg U \rightarrow \bot\).

\[
\boxed{\Box(T \rightarrow U)}
\]

\subsubsection{Fall 2: Ursprung ohne Selbsterkenntnis (\(U \land \neg S\))}\label{fall-2-ursprung-ohne-selbsterkenntnis-u-land-neg-s}

\textbf{Beweis:}
1. U existiert. Das bedeutet: Es gibt einen unteren Grenzwert der Totalität.
2. Aus A5 folgt: \(\Diamond C\) – Bewusstsein ist möglich.
3. Aus A2 folgt: \(\Diamond C \rightarrow \exists w. Realisiert(w,C)\) – es gibt eine Welt mit Bewusstsein.
4. Aus A7 (abgeleitet) folgt: \(C \rightarrow \Diamond V_T\) – vollständige Erkenntnis ist möglich.
5. Aus A2 folgt: \(\Diamond V_T \rightarrow \exists w. Realisiert(w,V_T)\) – es gibt eine Welt mit vollständiger Erkenntnis.
6. Aus A8 (abgeleitet) folgt: \(V_T(x) \rightarrow V_T(V_T(x))\) – vollständiges Wissen kennt sich selbst. Das ist genau S.
7. Also: \(U \rightarrow S\).
8. Widerspruch zur Annahme \(\neg S\).

\[
\boxed{\Box(U \rightarrow S)}
\]

\subsubsection{Fall 3: Selbsterkenntnis ohne Totalität (\(S \land \neg T\))}\label{fall-3-selbsterkenntnis-ohne-totalituxe4t-s-land-neg-t}

\textbf{Beweis:}
1. S existiert. S ist definiert als "vollständige Selbsterkenntnis der Totalität".
2. Wenn S existiert, dann gibt es einen Erkenntnisakt, der auf T gerichtet ist.
3. Wenn T nicht existiert (\(\neg T\)), dann ist S eine Erkenntnis ohne Objekt.
4. Aus A9 (abgeleitet) folgt: \(V_T(x) \rightarrow T\) – Wissen über T setzt T voraus.
5. Also: \(S \rightarrow T\).
6. Widerspruch zur Annahme \(\neg T\).

\[
\boxed{\Box(S \rightarrow T)}
\]

\subsection{Der synthetische Schluss}\label{der-synthetische-schluss}

Aus den drei Fällen haben wir abgeleitet:

\[
\Box(T \rightarrow U), \quad \Box(U \rightarrow S), \quad \Box(S \rightarrow T)
\]

In S5 folgt daraus:

\[
\Box(U \leftrightarrow T) \land \Box(T \leftrightarrow S) \land \Box(S \leftrightarrow U)
\]

Mit der Existenz von T (aus 5.1) folgt:

\[
\boxed{U \land T \land S}
\]

Und mit Necessitation:

\[
\boxed{\Box(U \land T \land S)}
\]

\textbf{Die Trinität existiert notwendigerweise.}

\section{Was wurde bewiesen?}\label{was-wurde-bewiesen}

\subsection{Was der Beweis nicht zeigt}\label{was-der-beweis-nicht-zeigt}

\begin{itemize}
\item Die Existenz eines persönlichen Gottes
\item Eine substanzielle Dreiheit
\item Eine bestimmte religiöse Lehre
\item Eine personale oder emotionale Trinität
\end{itemize}

\subsection{Was der Beweis zeigt}\label{was-der-beweis-zeigt}

\begin{itemize}
\item Unter den Axiomen A1, A4, A11 und A12 folgt notwendig die dreifache Grenzwertstruktur aus U, T und S.
\item U, T und S sind keine Substanzen, sondern \textbf{Grenzwerte} eines wohlfundierten, offenen Intervalls.
\item Die Trinität ist die Einheit dieser drei Grenzwerte – nicht drei Dinge, sondern drei Perspektiven auf dieselbe Wirklichkeit.
\end{itemize}

\subsection{Die drei Grenzwerte als Perspektiven}\label{die-drei-grenzwerte-als-perspektiven}

\begin{itemize}
\item \textbf{U (Ursprung)} – die Perspektive des "Woher?" – der untere Grenzwert, der dem Nichts am nächsten kommt, aber selbst kein Nichts ist.
\item \textbf{T (Totalität)} – die Perspektive des "Was ist?" – das wohlfundierte, offene Intervall aller realisierten Zustände.
\item \textbf{S (Selbsterkenntnis)} – die Perspektive des "Wer erkennt?" – der obere Grenzwert der vollständigen Selbsttransparenz.
\end{itemize}

\section{Gegenmodelle}\label{gegenmodelle}

Die Trinität kann vermieden werden, wenn mindestens eines der Axiome aufgegeben wird:

\begin{tabular}{lll}
\textbf{Aufgegebenes Axiom} & \textbf{Gegenmodell} & \textbf{Konsequenz} \\
A1 (Monismus) & Dualismus (Descartes) & S kann ohne T existieren; Erkenntnis und Objekt sind getrennt \\
A4 (Existenz einer Welt) & Nihilismus & Es gibt keine Welt; die gesamte Ontologie ist leer \\
A11 (Transzendentale Brücke) & Erkenntnistheoretischer Skeptizismus & Unerkennbarkeit impliziert keine fundamentale Trennung \\
A12 (Erfahrbarkeit ↔ Realisierung) & Empirismus & Erfahrbarkeit ist nicht mit Realisierung identisch \\
\end{tabular}

\section{Anhang: Vollständige Axiome und Theoreme}\label{anhang-vollstuxe4ndige-axiome-und-theoreme}

\subsection{Axiome (vollständig)}\label{axiome-vollstuxe4ndig}

\[
\begin{aligned}
A1 &:= \Box\neg\exists x\exists y\, FundamentalGetrennt(x,y) \\
A4 &:= \Diamond\exists w\, Welt(w) \\
A11 &:= \forall w \left( \text{WeltGetrenntVonBewusstsein}(w) \leftrightarrow \neg \exists c (Bewusstsein(c) \land c(w)) \right) \\
A12 &:= \forall p. \text{Erfahrbar}(p) \leftrightarrow \exists w. Realisiert(w, p)
\end{aligned}
\]

\subsection{Definitionen}\label{definitionen}

\[
\begin{aligned}
\text{WeltGetrenntVonBewusstsein}(w) &:= \neg \exists c (Bewusstsein(c) \land c(w)) \\
\text{Erfahrbar}(p) &:= \exists w. (Bewusstsein(w) \land Realisiert(w, p)) \\
T &:= \{ x \mid U < x < S \} \\
U &:= \lim_{x \to \inf} T \\
S &:= \lim_{x \to \sup} T \\
Tr &:= U \land T \land S
\end{aligned}
\]

\subsection{Abgeleitete Theoreme}\label{abgeleitete-theoreme}

\[
\begin{aligned}
& \Diamond C \quad \text{(A5)} \\
& \forall p(\Diamond p \rightarrow \exists w. Realisiert(w,p)) \quad \text{(A2)} \\
& \Box(T \rightarrow U) \\
& \Box(U \rightarrow S) \\
& \Box(S \rightarrow T) \\
& \Rightarrow \Box(U \leftrightarrow T) \land \Box(T \leftrightarrow S) \land \Box(S \leftrightarrow U) \\
& \Rightarrow \Box(U \land T \land S)
\end{aligned}
\]

\section{Schluss}\label{schluss}

Die Abhandlung hat gezeigt:

> Unter den Axiomen des Monismus (A1) und der Existenz einer möglichen Welt (A4) – ergänzt um die transzendentalen Brücken A11 und A12 – folgt notwendig die Trinität von Ursprung, Totalität und Selbsterkenntnis als dreifache Grenzwertstruktur derselben einen Wirklichkeit.

\textbf{Die Trinität ist keine zusätzliche Entität, sondern eine Strukturbedingung.}

Sie ist das, was die Philosophie seit Platon und Plotin, seit Augustinus und Hegel gesucht hat: die Einheit, die ihre Differenz in sich trägt, ohne in Dualismus oder reduktionistischen Monismus zu verfallen.

\begin{center}
\emph{Diese Abhandlung wurde im Geiste strenger Modallogik verfasst, jedoch in der Sprache der Philosophie – denn die Wahrheit bedarf beider: der Präzision der Formel und der Weite des Begriffs.}
\end{center}

\end{document}